% SPDX-License-Identifier: CC-BY-4.0
% Copyright (c) 2026 Martin Schröder <info@swedishembedded.com>

\section{Discussion and limitations}

\paragraph{Alternative explanations a gain must survive.} A gain on a repository's
reserved families can come from (a) learning the repository, (b) learning the tool
protocol of the harness, (c) learning to undo the synthetic mutation operators rather
than to debug, or (d) memorisation of public code. E3's off-diagonal cells and A1 address
(a) against (b) and the value of verification; (c) is addressed only by adding
historical bug-fix tasks to the exam, which the mutation-seeded generator does not produce;
(d) is bounded, not removed, by comparing with the base on the same exam.

\paragraph{Limits.} One model size (a Qwen3 base of at most 8B parameters, set by what
trains and serves end to end on two 24~GB GPUs), one agent harness, Python fixtures first
with Rust repositories later because offline dependency vendoring needs a working
container backend, runs without process isolation in this environment, and exam sizes of
a single repository that may be too small for a significant sign test. Agent rollouts at
the lengths of real trajectories make training cost unmeasured for the 8B model until
the long-context cost is measured.

\paragraph{If the thesis fails.} A null on E2 with a working E1 means the pipeline runs and
does not help at these sample sizes and this data volume; it would be reported as
such, with the discordant counts, not as a defect of the idea.

\section{Conclusion}

We fixed a protocol for asking whether a local agent can be tuned to a repository,
repeatedly and cumulatively, without a person in the loop, stated what its primary test can
decide, and listed what had to be repaired in the stack first. We have not yet measured
whether it works.
